\section{The stack, its lifts, and the nonzero instanton count}
\label{sec:stack}

We now arrange the ordinary maps and the fixed cap pairing into a closed
family with a nonzero integer count.  The cap pairing is nonzero over
\(\Q\), while the maps are units modulo two.  Repetition on the free
integral Floer lattice will reconcile these two inputs and make the
coupled Dirac index positive.  Before choosing that repetition, we compute
the stack's lift data and identify its actual count in the metrics used
for the spinor argument.  That identification requires a separate
exclusion of central limits at rational seams; invariance of the
individual maps alone does not supply it.

Throughout this section, the diffeomorphism
\(\phi:Y_{-2}\longrightarrow Y_2\) is the arbitrary oriented cosmetic
identification fixed in Section~\ref{sec:ordinary}.  No compatibility with
the knot exterior is imposed.  We use the caps of
Theorem~\ref{thm:caps:pairing}, their fixed compact probes of total real
degree \(z\), and the signed negative interval map
\(B:C(Y_2)\longrightarrow C(Y_{-2})\) of
Theorem~\ref{thm:negative:unit}.  Put
\[
 H=f_+f_-:C(Y_{-2})\longrightarrow C(Y_2).
\]
Thus \(H\) is the ordinary rigid map along
\(-2\to-1\to0\to1\to2\), with the odd determinant at \(Y_0\).
All compositions below use fixed determinant transitions and
determinant-one gauges.

\subsection{The arrangement and its cell lattices}

Insert \(n\) negative intervals \(W'_1,\ldots,W'_n\) between the caps.
After an interval use \(\phi\) to return to the \(Y_2\) convention,
except at \(m\) selected bridges, where the positive path defining
\(H\) is inserted instead.  The return after the last interval meets
the final cap in the \(Y_2\) convention.  Selected bridges are at least
four intervals apart; the repeated pattern used below has exact spacing
four.  Long initial and terminal pads contain no positive bridge.
A positive bridge together with its two
adjacent negative intervals is a \emph{macro}.  Macros are disjoint and
have ordinary single intervals between them.  Write \(J_i\) for the
potential primary \(S^3\)-cut in \(W'_i\), and write
\[
 S_i=F_{l,i}-F_{r,i},\qquad S_i^2=-4,
\]
for its test sphere.  Cutting at all the \(J_i\)'s and filling the new
spherical boundaries by balls gives two outside pieces and one cell
between each pair of successive \(J_i\)'s.  This is a topological
decomposition used for intersection-form computations; it does not
assert that any metric neck has been stretched.

\begin{figure}[ht]
\centering
\begin{tikzpicture}[
 block/.style={draw,minimum width=9mm,minimum height=8mm,inner sep=2pt},
 arrow/.style={-{Stealth[length=1.6mm]}},font=\small]
 \node (left) at (0,0) {$C_-$};
 \node[block] (b1) at (1.35,0) {$B$};
 \node[block,fill=blue!7] (h) at (2.7,0) {$H$};
 \node[block] (b2) at (4.05,0) {$B$};
 \node[block] (b3) at (5.65,0) {$B$};
 \node[block] (b4) at (7.25,0) {$B$};
 \node (dots) at (8.85,0) {$\cdots$};
 \node (right) at (10.2,0) {$C_+$};
 \draw[arrow] (left)--(b1);
 \draw[arrow] (b1)--(h);
 \draw[arrow] (h)--(b2);
 \draw[arrow] (b2)--node[below] {$\phi$}(b3);
 \draw[arrow] (b3)--node[below] {$\phi$}(b4);
 \draw[arrow] (b4)--node[below] {$\phi$}(dots);
 \draw[arrow] (dots)--(right);
 \draw[decorate,decoration={brace,amplitude=4pt}]
   (0.85,.65)--(4.55,.65)
   node[midway,above=5pt] {one positive macro};
 \draw[densely dashed] (1.35,-.4)--(1.35,-.85)
   node[below,font=\scriptsize] {$J_j$};
 \draw[densely dashed] (4.05,-.4)--(4.05,-.85)
   node[below,font=\scriptsize] {$J_{j+1}$};
 \draw[decorate,decoration={brace,amplitude=4pt,mirror}]
   (1.35,-1.35)--(4.05,-1.35)
   node[midway,below=5pt,font=\scriptsize] {positive cell $P$};
 \draw[decorate,decoration={brace,amplitude=4pt,mirror}]
   (0.85,-2.05)--(7.7,-2.05)
   node[midway,below=5pt] {four interval blocks per positive bridge};
\end{tikzpicture}

\caption{One repeated four-interval pattern between the caps.  The
first two intervals and the positive path form a macro; the positive
cell lies between the primary cuts inside those intervals.  The cut
marks are schematic, and the cell lattice is computed after filling
its spherical boundaries.  Initial and terminal negative pads are
suppressed.  All ordinary \(Y\)-seams have finite length; each
\(B\) carries its own primary \(J\)-face and lens face, which are
the genuine ends of its parameter interval.}
\label{fig:stack:pattern}
\end{figure}

\begin{proposition}[Cell lattices]
\label{prop:stack:lattice}
All these filled pieces have \(H_1=0\), and their integral intersection
lattices are unimodular.  A negative cell has orthogonal trace classes
\(F_r,F'_l\), both of square \(-2\), and its lattice has the basis
\begin{equation}
 a=\frac{F_r+F'_l}{2},\qquad b=\frac{F_r-F'_l}{2},\qquad
 a^2=b^2=-1,\quad ab=0.
\label{eq:stack:negative-lattice}
\end{equation}
A positive cell \(P\) has an integral basis
\(G,U,T_1,T_2,T_3,T_4\) with
\begin{equation}
 Q_P=\begin{pmatrix}2&1\\1&0\end{pmatrix}\oplus(-I_4),\qquad
 F_r=G-\sum_{b=1}^4T_b,\qquad F'_l=G-2U.
\label{eq:stack:positive-lattice}
\end{equation}
Here \(U\) is represented by a sphere of square zero, and
\(UF_r=UF'_l=1\).  The middle \(0\)-trace class is
\begin{equation}
 F_0=G-T_3-T_4.
\label{eq:stack:middle-class}
\end{equation}
The macro contains an ordinary plumbing of the three spheres
\begin{equation}
 X_1=-S_j,\quad U,\quad X_2=S_{j+1},\qquad
 X_1^2=X_2^2=-4,\quad UX_1=UX_2=1,\quad X_1X_2=0.
\label{eq:stack:macro-plumbing}
\end{equation}
In particular \(b^+(P)=1\) and \(\sigma(P)=-4\).
\end{proposition}

\begin{proof}
The trace halves have zero first homology.  Gluing two such halves along
a rational homology sphere preserves this property by the
Mayer--Vietoris sequence.  The same argument applies to the caps and to
the positive handle path.  Filling spherical boundary components by
balls does not change \(H_1\).  Poincar\'e duality and the universal
coefficient theorem therefore give a unimodular intersection form on
the integral second homology.

In a negative cell the two trace classes have intersection matrix
\(\operatorname{diag}(-2,-2)\), whose determinant is four.  They span
an index-two sublattice.  Its discriminant group is
\((\Z/2)^2\).  Of its three nonzero elements, the classes of
\(F_r/2\) and \(F'_l/2\) have square \(-1/2\), and so cannot be
adjoined to an integral lattice.  The remaining element,
\((F_r+F'_l)/2\), has integral square \(-1\).  It gives the unique
integral index-two extension, namely \eqref{eq:stack:negative-lattice}.
This argument uses no information about \(\phi\) beyond its being a
diffeomorphism of the filled manifolds.  In particular, its action on
\(H_1(Y_{\pm2})=\Z/2\) is necessarily the identity.

For the positive cell use the polygon model from
Section~\ref{sec:ordinary}.  Between its two spherical boundaries the
consecutive rays are
\[
 (-1,0),\ (-2,1),\ (-1,1),\ (0,1),\ (1,1),\ (2,1),\ (1,0).
\]
The five internal side spheres form the chain
\([-1,-2,-2,-2,-1]\).  The sum of these spheres has zero intersection
with each member of the chain.  Smoothing its four transverse
intersections gives a sphere: the dual graph is a tree, and each
smoothing joins two spheres.  Its square is
\(-1-2-2-2-1+2\cdot4=0\).  This is \(U\).
Its two endpoint incidences give the stated intersections with the
neighboring test spheres.

Blow down the four corner subdivisions introduced by the four
increasing steps.  The remaining doubled \(2\)-trace has a trace class
\(G\) with \(G^2=2\).  Doubling its cocore disk across the
\(Y_2\) seam gives the sphere \(U\), with \(U^2=0\) and
\(GU=1\).  Under restoration of the corner subdivisions this is
the smoothed side-chain sphere described above; the blowup points
are disjoint from its representative.  This rank-two form is unimodular.
Restore the subdivisions in their chronological order from slope
\(-2\) to slope \(2\), and use their total exceptional transforms
as \(T_1,\ldots,T_4\).  These have intersection form \(-I_4\) and
are orthogonal to the pullbacks of \(G,U\).  The incoming trace passes
through all four blown-up corners, whereas the middle trace passes
through the last two.  Their proper transforms are therefore
\(G-\sum T_b\) and \(G-T_3-T_4\).  The reverse outgoing trace is
orthogonal to the incoming one, meets \(U\) once, and has square
\(-2\); in the doubled-trace coordinates it is \(G-2U\).
This proves \eqref{eq:stack:positive-lattice} and
\eqref{eq:stack:middle-class}.  The local corner charts give transverse
positive plumbing intersections, so the smoothing and the two adjacent
\(-4\)-spheres give \eqref{eq:stack:macro-plumbing} inside the macro.
There is no use of a negative bridge identification in this construction.
The signature and positive rank follow from the displayed matrix.
\end{proof}

An evaluation vector \(v=(x,u;t_1,t_2,t_3,t_4)\) in the positive
basis consequently has square
\begin{equation}
 v^2=2xu-2u^2-\sum_{b=1}^4t_b^2.
\label{eq:stack:positive-dual-square}
\end{equation}
On a negative cell the corresponding formula is
\begin{equation}
 v^2=-\frac{(vF_r)^2+(vF'_l)^2}{2}.
\label{eq:stack:negative-dual-square}
\end{equation}
These formulas will also be used on rational orthogonal summands of
pieces cut at only some of the \(J_i\)'s.

\subsection{Characteristic lifts and the index telescope}

The spin-\(c\) structure used for the coupled Dirac operator will have
characteristic determinant class \(l\).  For a rank-two bundle with
determinant class \(c\), write \(\Lambda=l+c\).  We choose \(c_0\)
and \(\Lambda_0\) so that \(l=\Lambda_0-c_0\) is characteristic
and remains fixed when the determinant lift changes.

On every negative cell set
\[
 c_0=0,\qquad (\Lambda_0F_r,\Lambda_0F'_l)=(-2,0).
\]
On a positive cell specify their evaluations by
\begin{equation}
 c_0=(2,1;1,0,1,0),\qquad
 \Lambda_0=(-2,-1;0,1,0,-1).
\label{eq:stack:base-lifts}
\end{equation}
On each outside piece use the fixed cap determinant, extending it with
zero evaluation on the bordering negative trace.  Choose
\(\Lambda_0-c_0\) characteristic, with
\(\Lambda_0F_{l,1}=0\) on the first outside trace and
\(\Lambda_0F_{r,n}=-2\) on the last.  On the separate extra exceptional
class \(E_{\rm exc}\) of each cap choose
\(\Lambda_0E_{\rm exc}\) to be an odd integer; its absolute value
will later be made large.  The old cap data are fixed before this choice.

\begin{lemma}[Existence and restrictions of the lifts]
\label{lem:stack:lifts}
These choices define integral classes on the closed stack, and
\(l=\Lambda_0-c_0\) is characteristic.  For every bit vector
\(e=(e_1,\ldots,e_n)\in\{0,1\}^n\), set
\begin{equation}
 c_e=c_0+\sum_{i=1}^ne_i\PD(S_i),\qquad
 \Lambda_e=\Lambda_0+\sum_{i=1}^ne_i\PD(S_i).
\label{eq:stack:bit-lifts}
\end{equation}
Then \(\Lambda_e-c_e=l\) is independent of \(e\).  The determinant
has the required odd restriction at every positive \(Y_0\) seam and
is trivializable at all rational primary seams.  Moreover
\begin{equation}
 c_0S_i=0,\qquad \Lambda_0S_i=2,\qquad
 S_iS_j=0\quad(i\ne j).
\label{eq:stack:sphere-pairings}
\end{equation}
\end{lemma}

\begin{proof}
In the basis \(a,b\) of a negative cell, \(\Lambda_0\) has
evaluations \((-1,-1)\), so its difference from \(c_0=0\) is
characteristic.  On a positive cell this difference has evaluations
\[
 (-4,-2;-1,1,-1,-1).
\]
The first two are even, as the diagonal entries of the rank-two form
are even, and the other four are odd.  This is exactly the
characteristic congruence on the integral basis in
Proposition~\ref{prop:stack:lattice}.  Direct evaluation gives, for
either type of cell,
\[
 (c_0F_r,c_0F'_l)=(0,0),\qquad
 (\Lambda_0F_r,\Lambda_0F'_l)=(-2,0).
\]

A square-\(-2\) integral class in an integral lattice is primitive:
if it were \(dF\), \(d>1\), its square would be divisible by
\(d^2\).  In a unimodular lattice it consequently has a dual class
of pairing one.  Starting with any characteristic class on an outside
piece, add an even multiple of that dual to give the prescribed even
evaluation on the bordering trace.  The extra exceptional summand is
orthogonal and has determinant evaluation zero; its characteristic
evaluation can be any odd integer.  Gluing over \(S^3\)'s identifies
these integral classes, since the first and second cohomology of the
seams vanish.  This proves existence of the closed classes and the
characteristic assertion.

On a positive cell \(c_0F_0=1\), and
\(F_0F_r=F_0F'_l=0\).  Thus every adjacent bit change preserves the
odd evaluation at the middle seam.  At a \(\pm2\) rational seam the
restriction of a line class is its trace evaluation modulo two.
Those evaluations are even, including after a bit change.  Hence the
determinant is trivializable there.  The remaining rational seams have
first homology zero.  Finally the two halves of each \(S_i\) have the
base evaluations just computed, including at the outside pieces.
The spheres lie in distinct intervals, which proves
\eqref{eq:stack:sphere-pairings}.
\end{proof}

Write
\(\kappa=c_2(E)-c_e^2/4\) and
\(\Theta=(\Lambda_e^2-\sigma(X))/4\).  To keep \(\kappa\) fixed
when one bit changes from zero to one, the second Chern number must
\emph{decrease} by one.  Indeed
\begin{equation}
 c_e^2=c_0^2-4\sum_i e_i,\qquad
 \Lambda_e^2=\Lambda_0^2.
\label{eq:stack:global-bit-squares}
\end{equation}
There is a useful more local version of the second equality.

\begin{lemma}[The cell telescope]
\label{lem:stack:telescope}
The contribution of the cell between \(J_j\) and \(J_{j+1}\) to
\(\Theta\) is
\begin{equation}
 m_{\rm cell}+\frac{e_j-e_{j+1}}2,
 \qquad
 m_{\rm cell}=\begin{cases}0&\text{negative cell},\\1&\text{positive cell}.
 \end{cases}
\label{eq:stack:cell-theta}
\end{equation}
Consequently there are constants \(b_0,\Theta_0\), depending only
on the outside choices, such that
\begin{equation}
 b^+(X)=m+b_0,\qquad \Theta=m+\Theta_0.
\label{eq:stack:total-topology}
\end{equation}
For an interior filled component consisting of \(L\) cells between
actual \(J\)-cuts, with \(m_I\) positive cells and endpoint bits
\(e_a,e_b\),
\begin{equation}
 b_I^+=m_I,\qquad \Theta_I=m_I+\frac{e_a-e_b}{2},\qquad p_I=L-1.
\label{eq:stack:interior-topology}
\end{equation}
Here \(p_I\) counts interval parameters with separated lens parameters
virtually restored.  Endpoint differences telescope when contiguous
components are joined.
\end{lemma}

\begin{proof}
The restriction of a bit change to a cell is
\(-e_j\PD(F_r)+e_{j+1}\PD(F'_l)\).  Orthogonality of the two
trace classes and their base evaluations give
\[
 \Lambda_e^2-\Lambda_0^2
 =4e_j-2e_j^2-2e_{j+1}^2=2(e_j-e_{j+1}).
\]
The base square on a negative cell is \(-2\), equal to its
signature.  On a positive cell both base squares in
\eqref{eq:stack:base-lifts} are zero by
\eqref{eq:stack:positive-dual-square}, whereas its signature is
\(-4\).  This proves \eqref{eq:stack:cell-theta}.
Summing the internal cells leaves \((e_1-e_n)/2\).
The first outside trace contributes \(-e_1/2\), and the last
contributes \(+e_n/2\).  All bit dependence cancels.  Positive
subspaces add across spherical cuts, proving
\eqref{eq:stack:total-topology}.  The same sum on an interior segment
proves \eqref{eq:stack:interior-topology}; there is one retained
interval between each consecutive pair of its \(L\) cells.
\end{proof}

\begin{corollary}[Spacing]
\label{cor:stack:interior}
An interior segment with \(m_I\ge1\) positive cells has
\(L\ge4m_I-3\).  More precisely, if \(a,b\) count negative cells
before its first and after its last positive cell, then
\(L\ge4m_I-3+a+b\).  Equality holds when it contains no enlarged
gap between positive cells.
\end{corollary}
\begin{proof}
The positive cell positions are at least four apart.  The distance
from the first to the last is therefore at least \(4(m_I-1)\).
Include those endpoint cells and the \(a+b\) additional cells.
\end{proof}

The dimension of the ordinary closed family with its \(n\) sphere
tests and cap probes is \(D_I+n-2n-z\).  Thus the zero-dimensional
condition is
\begin{equation}
 D_I=8\kappa-3(1+b^+)=n+z.
\label{eq:stack:ordinary-dimension}
\end{equation}
The coupled Dirac index is \(n_D=\Theta-\kappa\), as computed in
Lemma~\ref{lem:indices:closed}.  Consequently
\begin{equation}
 n_D=\Theta-\kappa
 =\frac{5m-n}{8}+\Theta_0-\frac{z+3+3b_0}{8}.
\label{eq:stack:dirac-growth}
\end{equation}
These are topological index formulas; none presupposes the existence
of a solution.  For the four-interval repeated pattern with fixed pads,
\(n=4m+d\) for a fixed integer \(d\), so the last index grows as
\(m/8\).  We still have to identify the actual ordinary count with
the product and choose repetitions that preserve the cap pairing and
the bundle-degree congruence.

\subsection{Tests and the relative macro map}

Use a product of the negative interval parameters.  The genuine faces
are the primary \(J\)-faces and the \(-4\)-lens faces of those
intervals.  All other seams, including the outer \(Y\)-seams of
macros and single intervals, have long \emph{finite} collars.
The metric construction of Theorem~\ref{thm:geometry:family} keeps
these genuine faces and provides the local metrics later needed for
the Abelian estimate.  In the present ordinary argument we use only
the topology, the face-product structure, and the small-perturbation
properties stated next; no Abelian spinor inequality is used.

\begin{lemma}[Compatible sphere representatives]
\label{lem:stack:tests}
The sphere tests can be chosen with the following properties, preserving
the integral map \(B\) on homology and its mod-two unit property.
For a commuting connection whose ordered line difference is \(v\),
a whole test avoids \(-1\) if \(vS_i=0\).  Toward the primary
face it becomes two consecutive sweeps representing
\(F_{l,i}\) and \(-F_{r,i}\); each sweep then avoids \(-1\)
when its half evaluation is zero.  During the change of prescription,
avoidance is retained if both half evaluations vanish.  At an actual
\(J_i\)-split each surviving test is assigned compactly to exactly
one side.  At a lens face it is supported on the lens cap and its data
at the two determinant lifts are tensor-paired as in
Proposition~\ref{prop:negative:local-sign}.
\end{lemma}

\begin{proof}
Apply the smooth conjugation-equivariant operation of
Lemma~\ref{lem:analysis:zero-winding} to the normalized holonomy path
of each sphere sweep.  At a commuting connection its winding is half
the line-difference evaluation.  The operation fixes nonzero-winding
commuting paths and moves every zero-winding commuting path into the
closed positive-real hemisphere, giving a uniform gap from \(-1\).
It is defined on path space before evaluating on connections, so one
smooth prescription applies across all commuting configurations.

At a primary split the simply connected traces supply sphere maps for
\(F_{l,i}\) and \(-F_{r,i}\), and their pinch represents \(S_i\).
Lemma~\ref{lem:analysis:segmentation} first homotopes the whole sweep
to this concatenation and then interpolates to the separate modifications
of the two half sweeps.  When both half evaluations vanish, the whole
and half logarithms agree on their respective segments, so the uniform
avoidance gap persists throughout this interpolation.  At the actual
face each event is defined in a frame intrinsic to its half; a zero
half evaluation therefore gives avoidance on that side.  These changes
take place within the interval, before its true length collars and away
from the internal \(Y\)-necks used below.  The metric construction
stops using the whole prescription on a side before that side begins
to switch.

Use product data near primary faces and the original tensor-paired
data near lens faces.  The two representative lemmas belong to the
independent ordinary library in
Proposition~\ref{prop:analysis:ordinary-library}; they do not use the
nonzero count proved below.  Generic equivariant cut variations smaller
than half the uniform gap retain the avoidances and give regular cuts
on the remaining irreducible locus.

To verify invariance of \(B\), use the homotopy to the path operation
provided by Lemma~\ref{lem:analysis:zero-winding}.  It is relative to
the mark endpoints and fixes the constant path and every nonzero-winding
commuting path, including the minimum lens-cap paths of
Lemma~\ref{lem:negative:cap}.  Thus flat avoidance and tensor
compatibility at the lens are retained.  The operation uses only the
original sweep; its exact-support and ideal-persistence properties are
those of Lemma~\ref{lem:analysis:incidence}, and equation data remain
independent of the incidence mark.  This initial homotopy need not avoid
\(-1\) at every intermediate winding-zero path, nor be a small change
of sweep: Proposition~\ref{prop:negative:invariance} imposes neither
condition.  The later whole-to-halves interpolation has the stronger
both-zero-half avoidance just proved.  At a primary face the
unconstrained index exclusion applies even without the test, and at a
lens face the two lift contributions cancel by
Proposition~\ref{prop:negative:local-sign}.  The invariance proposition
therefore preserves \(B\) on homology throughout these changes,
including for standalone intervals.
\end{proof}

In estimates involving ideal limits, failure of an avoidance can occur
only when a particle meets the exact support of that test.  Distinct
tests have disjoint supports in their intervals.  One lost unit of
charge cannot be used to explain two such failures.  The ordinary
transversality and support-incidence estimates are those of
Lemma~\ref{lem:ordinary:perturbations}.

\begin{lemma}[Local charge costs]
\label{lem:stack:charge-costs}
For the main pieces of a single interval or macro, virtually restore
all separated lens caps with the split references
of differences \(0,\pm4,\pm2\), respectively at the extendibly
central, other central, and trace-zero states.  Include their charges,
particles, and other positive losses in the piece charge.  For a
reducible main piece with \(s\) assigned tests and no positive cell,
\begin{equation}
 8\kappa_{\rm piece}\ge4s.
\label{eq:stack:negative-test-cost}
\end{equation}
For a reducible main piece containing the positive cell of a macro,
\begin{equation}
 \kappa_{\rm piece}\in\tfrac12\Z,\qquad
 \kappa_{\rm piece}\ge\tfrac12.
\label{eq:stack:positive-charge}
\end{equation}
The statements hold for sufficiently small perturbations, throughout
the finite metric templates and the long-neck comparisons used below.
\end{lemma}

\begin{proof}
For an ordered splitting the line difference satisfies \(v=c_e\)
modulo two.  All negative trace evaluations are therefore even.
After virtual filling its charge has the form
\[
 \kappa_{\rm piece}=-v^2/4+\ell,\qquad
 \ell\in\Z_{\ge0}.
\]
The stated reference fillings have charges \(0,1,1/4\); their
boundary eigenlines extend the given splitting.  Excess charge in a
cap or its trajectories, and charge of particles, contributes to
\(\ell\).  On a negative piece a nonzero even half evaluation
contributes at least two to \(-v^2\), by
\eqref{eq:stack:negative-dual-square}.  Every assigned test either
charges such a half or, by Lemma~\ref{lem:stack:tests}, charges a
unit of \(\ell\).  In a whole test a nonzero difference entails a
nonzero half; during switching the same assertion follows unless
both halves vanish; at a primary face only its assigned half is used.
For a separated lens, vanishing of both halves would give the flat
reference, which misses the test.  Different tests use different
halves or different units of loss.  Thus each costs at least four in
\(8\kappa=-2v^2+8\ell\), proving
\eqref{eq:stack:negative-test-cost}.

On the positive cell, \(vG\) is even and precisely two of the four
\(vT_b\)'s are odd; bit changes preserve these two properties.  Equation~
\eqref{eq:stack:positive-dual-square} shows that \(v^2\) is even.
Added negative trace halves also have even square contribution, so
\(\kappa\in\frac12\Z\).  The small-perturbation energy bound
first gives \(\kappa\ge0\) in this discrete set.  Charge zero is
impossible.  In fact, the sphere between the first two increasing
steps represents \(\pm(T_2-T_1)\), and the determinant evaluates
oddly on it.  It lies in the first half of the positive path, away
from the internal \(Y_{-2},Y_0,Y_2\) collars and from the genuine
cut collars.  For fixed artificial metric sizes its neighborhood has
compactly controlled geometry.  Were charge zero possible for
perturbations tending to zero, the energy identity would force the
trace-free curvature to tend to zero on this neighborhood.  A flat
projective limit over this sphere has trivial \(w_2\), contradicting
the odd determinant evaluation.

This argument remains uniform in the long-neck comparisons.  On a
long Floer collar retain a gradient perturbation; the perturbed
Chern--Simons identity bounds its charge below by an error tending
to zero with the perturbation size, independently of collar length.
The endpoint perturbation terms telescope over successive levels.
On genuine separating collars use unperturbed deep regions, and
make other errors arbitrarily small in integrated square norm.
The compact odd sphere is untouched.  Thus the same contradiction
excludes zero in true split limits and in the comparison to long
internal product collars.  Discreteness now gives
\eqref{eq:stack:positive-charge}.
\end{proof}

\begin{proposition}[The relative macro]
\label{prop:stack:macro}
The two-parameter, two-test family on the chronological cobordism
\(B\,H\,B\), with the above data and irreducible outer limits,
defines an integral map on Floer homology at instanton index \(i=2\).
It is independent of the metric interpolation preserving the genuine
faces, and agrees, up to the common orientation convention, with
\(BHB\).  In particular it induces a unit modulo two.
\end{proposition}

\begin{proof}
We give the low-dimensional exclusions needed for the chain and
homotopy identities.  Let \(t=0\) for the family and \(t=1\) for
a homotopy.  It suffices to consider \(i\le3-t\).  If \(k\) of
the two primary \(J\)'s are cut, \(k\in\{0,1,2\}\), let \(p\)
and \(s\) on a main piece denote its remaining interval parameters
and its assigned tests, with lens parameters virtually restored.
Then
\[
 \sum p=2-k,\qquad \sum s=2.
\]
Equation and cut data on a piece depend only on its retained
parameters.  A regular irreducible piece costs at least \(2s-p\).
At most one additional parameter is subtracted from the irreducible
aggregate in a homotopy.

A separated lens loses a parameter and removes its degree-two test
from the main problem, but has framed index shift
\(\delta=8E\ge2\).  Its net additional cost, measured against
\(2s-p\), is at least one.  A unit particle loses eight indices;
its position and any released cut degrees recover at most four by
the exact-support incidence convention.  These losses are therefore
strict.  A purely negative reducible piece has index at least
\(4s-3\), by Lemma~\ref{lem:stack:charge-costs}.  Matching at each
central contact adds three.

Consider first an unwanted centered contact or a maximal consecutive
reducible group not containing a positive-cell reducer.  A group of
\(r\) negative reducers contributes at least \(4s_R-3r\), and
its \(r-1\) internal and two exterior central matchings contribute
\(3(r+1)\).  Its aggregate cost is at least \(4s_R+3\).
A centered contact without reducers has the same extra cost three.
Adding regular irreducible costs, and allowing the one homotopy
parameter, gives the lower bound
\[
 4s_R+3+2(2-s_R)-(2-k)-t
 \ge5+k-t>3-t.
\]
Unused parameters on reducible pieces only increase this bound.

Now suppose the unique positive cell lies on a reducible piece.
Its index is at least \(4-6=-2\).  A maximal reducible group
containing it, including all internal and both exterior central
matchings, costs at least four; negative members add their
nonnegative test costs.  At \(k=0\), the whole main cobordism is
reducible and must be joined to each prescribed irreducible outer
limit by a genuine nonconstant irreducible trajectory.  These two
trajectories give the bound six.  At \(k=1\), let \(u\in\{0,1\}\) record whether the split
test is assigned to the positive piece.  That piece reaches an outer
end, so a true-end trajectory is required.  The other piece has no
retained parameter and has \(1-u\) tests.  If it is irreducible,
the bound is \(4+1+2(1-u)-t\ge5-t\).  If it is reducible,
both outer ends require trajectories and the bound is
\(4+4(1-u)+2\ge6\).  At \(k=2\) there are no retained
interval parameters.  Let \(u_i\) record assignment of each split
test to the positive piece.  An irreducible outer piece contributes
at least \(2(1-u_i)\); a reducible outer piece, together with its
forced true-end trajectory, contributes at least \(4(1-u_i)+1\).
Including the positive group and subtracting the homotopy parameter
at most once gives \(4-t\).  Each bound is strictly greater than
\(3-t\).  These are all placements and assignment choices for
the positive reducer.

It follows that only regular irreducible pieces occur in the relevant
boundary strata.  Their strict loss estimates exclude particles and
all lens possibilities except a single trace-minimum lens cap in the
boundary index.  Its two lift contributions cancel by
Proposition~\ref{prop:negative:local-sign}.  Ordinary Floer breaking
and regular gluing consequently give the integral chain identity and
the invariance on homology.

It remains to identify this operation with \(BHB\).  Homotope the
local toric metrics to product
metrics through families with the same genuine cuts.  The space of
metrics on the remaining compact cores permits the interpolation;
there is no need to retain a spinor clamp during this homotopy.
Keep the tests of Lemma~\ref{lem:stack:tests} and the paired lens
data.  Finite partial stretches may still retain the positive
rank of the unsplit macro; they are covered by the preceding macro
homotopy bounds and Lemma~\ref{lem:stack:charge-costs}.  For the
final comparison send all three internal lengths to infinity so
that their minimum tends to infinity.  Every limit then contains
the three cuts \(Y_{-2},Y_0,Y_2\), and its main factors in path
order are
\[
 B_L,\qquad H_-:\ -2\to-1\to0,\qquad
 H_+:\ 0\to1\to2,\qquad B_R.
\]
Both \(B\) factors are negative definite.  The two increasing
paths also have \(b^+=0\).  Explicitly, preceding \(H_-\) by
the \(-2\)-trace gives \(C_0\#2\overline{\CP}^{2}\),
with basis \(A^2=0,T_1^2=T_2^2=-1\).  Its incoming class is
\(A-T_1-T_2\); the orthogonal image of the handle path has
\(t_1+t_2=0\) and square \(-2t_1^2\).  Preceding \(H_+\)
by the \(0\)-trace gives \(C_2\#2\overline{\CP}^{2}\),
with \(G^2=2,T_3^2=T_4^2=-1\) and incoming class
\(G-T_3-T_4\).  A vector \(aG+t_3T_3+t_4T_4\) in its
orthogonal satisfies \(2a+t_3+t_4=0\) and has square \(-(t_3-t_4)^2/2\).
The relative Mayer--Vietoris sequence identifies these orthogonal
images with the handle-path homology; the null direction is the
\(Y_0\) surface direction.  Thus neither path contributes a
positive subspace.

At a compatible genuine \(J\)-face, either \(B\) splits into
negative trace halves; at a lens face, virtual restoration of the
negative lens cap recovers the same nonpositive form.  The two
increasing paths remain unchanged.  A reducer cannot meet the
admissible \(Y_0\) seam, so all reducer groups lie in negative
interval pieces, trace halves, or rational cylinder levels.  They
therefore have the same cost \(4s_R+3\), including their internal
and two exterior central matchings.  Regular increasing-path
factors cost nonnegative index and carry no interval parameters or
tests.  Consequently an unwanted centered contact again costs at
least \(5+k\), while a lens, particle, or nonconstant irreducible
trajectory adds strictly to the regular index bound.  At \(i=2\)
only the ordinary rigid irreducible matches remain, with no true
interval faces or positive losses.  The regular gluing theorem therefore identifies the map with the product
for sufficiently long finite internal lengths.  The coherent
orientations are the compositions of the ordinary orientations, with
the same local flip comparisons.  The mod-two unit assertion follows
from Theorems~\ref{thm:ordinary:units} and
\ref{thm:negative:unit}.
\end{proof}

\subsection{The closed gluing budget}

We next stretch only the external negative \(Y\)-seams separating
caps, macros, and single intervals.  The relative macro proposition
fixed irreducible outer limits; a limit of this closed stretching may
instead have central flats on the seams.  We therefore need a local
index bound allowing central outer limits, followed by a global bound
including their matching contributions.

\begin{lemma}[Local excess, including central outer limits]
\label{lem:stack:local-excess}
For a single interval or macro let \(b=1\) or \(2\), respectively,
and let \(i_{\rm loc}\) be its index at its actual outer limits,
with seam trajectories treated separately.  A single interval has
\(i_{\rm loc}-1\ge0\).  A macro with \(a\) central outer limits
has
\begin{equation}
 i_{\rm loc}-2\ge-2a.
\label{eq:stack:macro-excess}
\end{equation}
If both outer limits are irreducible, equality for either kind of
factor requires an ordinary stable rigid solution: there are no
primary or lens faces, particles, or positive trajectory losses.
\end{lemma}

\begin{proof}
Suppose \(k\) primary cuts occur in the factor.  Restore the lens
caps virtually, assigning any primary-neck charges to the pieces.
Index addition and the parameter and test totals give
\begin{equation}
 E:=i_{\rm loc}-b
 =\sum_{\rm pieces}(i+p-2s)+4k,
 \qquad \sum p=b-k,\quad\sum s=b.
\label{eq:stack:excess-identity}
\end{equation}
Indeed primary matching adds \(3k\), while the lost parameters
supply the remaining \(k\).  Call the summands in parentheses
\emph{scores}.  Regular irreducibles have score at least zero,
strictly positive for the losses described in the preceding proof.
Negative reducers have score at least
\(-3+2s+p\).  A reducer containing the positive cell has score
at least \(-2+p-2s\).

For a single interval with \(k=0\), an irreducible has nonnegative
score, and a reducer has \(s=p=1\), also giving score at least zero.
For \(k=1\), both parameters disappear and the two test assignments
are \(u,1-u\), with \(u\in\{0,1\}\).  If both pieces are
irreducible, \(E\ge4\).  If just one is reducible, assigned
\(s_R\) tests, then \(E\ge1+2s_R\ge1\).  If both are
reducible, their scores plus four give \(E\ge0\); both outer
limits are then central.  This proves all assertions for a single
interval.

For a macro first suppose the positive-cell piece is irreducible.
At \(k=0\) its score is nonnegative.  At \(k=1\), the only
possible reducer is the negative outer piece; with \(u\) the
assignment of the split test to the positive piece, its score is
at least \(-1-2u\), so \(E\ge3-2u\ge1\).  If it is
irreducible as well, the stronger bound is four.  At \(k=2\),
let \(u_i\) indicate assignment of the \(i\)-th split test to
the positive piece and \(r_i\) indicate reducibility of the
corresponding outer piece.  Then
\[
 E\ge8-\sum_{i=1}^2r_i(1+2u_i)\ge2.
\]
Thus all primary faces are strict in this case.

It remains to consider a positive-cell reducer.  At \(k=0\) the
whole macro is reducible, with \(p=s=2\), giving \(E\ge-4\)
and \(a=2\).  At \(k=1\) its parameter count is one and its
test count is \(1+u\), where \(u\in\{0,1\}\) is the split
test assignment.  The other piece has no parameter and has
\(1-u\) tests.  Write \(r=1\) if that piece is reducible and
\(r=0\) otherwise.  Equation~\eqref{eq:stack:excess-identity}
becomes
\[
 E\ge1-2u-r(1+2u).
\]
For \((u,r)=(0,0),(1,0),(0,1),(1,1)\), the respective lower
bounds are \(1,-1,0,-4\).  The positive reducer forces one central
outer limit, and the other reducer, if present, forces the second.
Hence \(a\ge1+r\), and each of these four bounds implies
\eqref{eq:stack:macro-excess}.

At \(k=2\) the positive reducer has no parameter and has
\(u_1+u_2\) tests.  Each outer piece has no parameter and has
\(1-u_i\) tests.  Thus
\[
 E\ge6-2(u_1+u_2)-\sum_{i=1}^2r_i(1+2u_i).
\]
For \(r_1+r_2=0,1,2\), its minimum over the assignments is,
respectively, \(2,-1,-4\).  Since \(a\ge r_1+r_2\), all
three minima imply \eqref{eq:stack:macro-excess}.  Notice that
an irreducible outer piece can itself have a central outer limit;
only this inequality, not equality, is needed.  When \(a=0\)
the bound is at least two.

Finally, with irreducible outer limits a nonsplit whole reducer is
impossible.  The strict primary-face bounds just proved, and the
strict lens and particle costs on remaining irreducibles, show that
zero excess then occurs only in the ordinary stable rigid count.
\end{proof}

\begin{proposition}[Closed composition in the final metrics]
\label{prop:stack:composition}
Choose the old cap metrics with generic periods and take sufficiently
long finite isolation collars for their separate extra exceptional
summands.  For any fixed finite stack, the sphere tests and sufficiently
small ordinary data can be chosen so that, for all sufficiently long
finite external negative \(Y\)-necks, the oriented closed count
\(\Omega\), summed over the determinant bits, equals the ordinary
composition of its cap states, single intervals, and macro maps, up
to one overall sign.  The same conclusion holds after sufficiently
small further generic changes preserving the genuine-face
prescriptions.  These data can also satisfy
Theorem~\ref{thm:estimates:clamp}.
\end{proposition}

\begin{proof}
The fixed indices and total charge, together with uniformly bounded
end corrections, give bounded cap charge ranges for this finite stack.
Include also the cap-cycle, comparison-homotopy, and old closed-pairing
ranges.  With the old-cap periods fixed, Proposition~\ref{prop:caps:periods},
in the uniform isolation form of Proposition~\ref{prop:estimates:cap},
permits exceptional isolation depths and perturbation tolerances to be
chosen for all these ranges at once.  It excludes cap reducers also in broken
relative limits as the external \(Y\)-necks stretch.  The charge
bounds used there remain uniform in those lengths: the gradient
Chern--Simons terms telescope along any chain of end levels, and the
resulting cap and finite-prefix bounds suffice for the period test.
They make no assertion of finite unperturbed curvature norm on an
infinite tail approaching a perturbed nonflat.

Each regular cap irreducible now has nonnegative excess over its fixed
probe degree.  The last assertion of Theorem~\ref{thm:caps:pairing}
allows the cap pairing to be recomputed with these final data by
stretching the old closure's rational seam.  Its closed invariant is
unchanged, since \(b^+>1\) and the probes represent the same classes.
Thus the pairing stays nonzero without requiring separate metric
independence of the two cap states.

Now consider a limit as all external seams become long.  A seam
contributes nonnegative additional index from trajectories and
matching.  If any central state occurs in it, its total contribution
is at least three: group contiguous reducible cylinder levels and
include their two exterior matching terms, exactly as in
Lemma~\ref{lem:ordinary:central-cost}.  Write \(S\) for the number
of seams containing a central state, and \(A\) for the sum of the
numbers of central outer limits over the macros.  Since macros are
separated by single intervals and separated from the caps by pads,
a seam is adjacent to at most one macro.  Hence \(A\le S\).
Lemma~\ref{lem:stack:local-excess} and cap regularity give total
excess at least
\[
 3S-2A\ge S.
\]
But the total excess is zero by
\eqref{eq:stack:ordinary-dimension}.  Thus \(S=0\).
With all external limits irreducible, every factor has nonnegative
excess.  Equality and the last assertion of
Lemma~\ref{lem:stack:local-excess} exclude all true faces and
positive losses.  Nonconstant irreducible seam trajectories cost
at least one and are excluded as well.  Precisely the rigid
irreducible factor matches remain.

Choose regular ordinary data independently on the pieces, with
product compatibility at all external seams and genuine cuts.
Ordinary compactness, the preceding exclusions, and regular Floer
gluing now identify the closed count with the product for
sufficiently long finite external lengths.  This proves also that
no broken or ideal cut tuple occurs at a genuine face for these
choices.  Otherwise a sequence at arbitrarily large external
lengths would have a limiting tuple excluded by the same budget.
At fixed suitable lengths the absence is open by compactness, and
the regular zero-dimensional count is unchanged by sufficiently
small further data changes.

For orientations, order the closed ASD determinant line in
cobordism order, pairing an intermediate end line with its dual.
The parameter, mark, and test lines are then reordered into the
factor order.  Their dimensions and the indices of the rigid
factors are fixed, so the reordering sign is independent of the
matched solutions.  Equivalently, the fixed cap degree determines
the first Floer parity, and the fixed degrees of the successive
maps determine every subsequent parity.  This gives a single
comparison sign for the whole product.  The local two-lift
comparison commutes with this operation: excise its framed lens
cap before or after the other gluings, using the same transition
and trace reference.  The compared cap indices are equal and even.
Thus all bit comparisons are precisely the ones used to define the
maps in Section~\ref{sec:negative}.

It remains to explain compatibility with the later metric choices.
The metric existence assertion in
Theorem~\ref{thm:geometry:family} and the finite-choice assertion in
Corollary~\ref{cor:estimates:choices} apply to the fixed topology,
lifts, and index bounds, independently of the value of \(\Omega\).
They supply finite artificial isolation sizes and a positive
allowable size for fixed gradient perturbations on the external
\(Y\)-collars before those collars are sent into their gluing
range.  The ordinary argument above works for arbitrarily small
such gradients.  Choose them below that tolerance, then take the
external lengths sufficiently large and finite.  Any additional
non-gradient errors are chosen with arbitrarily small total squared
norm for those already chosen lengths.  The final open generic
adjustments preserve the count.  Thus no error estimate is required
to hold for a non-gradient perturbation accumulated over an
uncontrolled \(Y\)-length.
\end{proof}

Changing the fixed determinant connection, or the background connection
used in the coupled Dirac equation, does not change this ordinary
calculation.  In the ordinary projective equations its effect is
absorbed by the scalar half-difference of the two determinant
connections.  The normalized holonomy tests use the same determinant
lifts.  Thus the independent background choices needed later do not
alter \(\Omega\).

\subsection{Repetition on the integral free quotient}

The cap pairing is rationally nonzero, whereas the maps are known to
be units modulo two.  The passage between these assertions uses the
free quotient of integral Floer homology, rather than discarding
torsion inside a mod-two homology group.

\begin{lemma}[The finite-lattice argument]
\label{lem:stack:finite-lattice}
Let \(C,C'\) be finite complexes of finitely generated free abelian
groups, and let \(f:C\to C'\) be a chain map whose reduction modulo
two is a quasi-isomorphism.  On the free quotients
\[
 L=H(C;\Z)/\operatorname{Tor},\qquad
 L'=H(C';\Z)/\operatorname{Tor},
\]
its induced map is an isomorphism after localization at two.  In
particular \(L/2^NL\to L'/2^NL'\) is an isomorphism for every
\(N\ge1\).

If \(L=L'\) has rank \(r\), all these maps modulo \(2^N\)
have orders dividing a single integer depending only on \(r,N\).
\end{lemma}

\begin{proof}
The mapping cone of \(f\) is a finite free complex whose mod-two
homology is zero.  Apply the universal coefficient sequence over
\(\Z_{(2)}\).  In each degree the finitely generated localized
cone homology has zero quotient by its maximal ideal \((2)\),
so it vanishes by Nakayama's lemma.  Thus \(f\) induces an
isomorphism on integral homology after localization.  An isomorphism
preserves torsion, hence also induces an isomorphism on the free
quotients after localization.  In integral bases its square matrix
therefore has odd determinant.  The adjugate formula makes it
invertible modulo every \(2^N\).  The last assertion follows by
taking the exponent of the finite group
\(\GL_r(\Z/2^N\Z)\).  It does not require integral homology to
be torsion-free.
\end{proof}

\begin{theorem}[A nonzero family with positive coupled index]
\label{thm:stack:nonzero}
Fix the old caps and their compact probes.  Prescribe any lower bounds
on the two outside pad lengths, and prescribe any odd evaluations of
\(\Lambda\) on their separate extra exceptional classes.  There
exist arbitrarily large finite stacks with positive bridges at
spacing four, satisfying those pad bounds, for which the following
hold simultaneously:
\begin{enumerate}
\item the lifts are those of Lemma~\ref{lem:stack:lifts}, with a
common allowed charge \(\kappa\) and
\(D_I=n+z\);
\item the actual oriented ordinary family count satisfies
\(\Omega\ne0\);
\item the coupled Dirac index \(n_D=\Theta-\kappa\) is positive;
\item all auxiliary collars are finite, and the metric and small-data
choices can be made to satisfy Theorem~\ref{thm:geometry:family} and
Theorem~\ref{thm:estimates:clamp}.
\end{enumerate}
The integers specifying the stack and all topological lift data are
chosen before the artificial isolation sizes, perturbation bounds,
and external Floer gluing lengths.
\end{theorem}

\begin{proof}
Clear the fixed cap denominators once.  Thus choose integral cap
cycles and an integral cap functional whose pairing is a nonzero
integer \(q\); the corresponding exact surface insertions are the
fixed integral multiples of the rational probes used in
Theorem~\ref{thm:caps:pairing}.  The functional kills torsion, so
it factors through the free integral Floer lattice \(L\) of
\(Y_2\).  Choose \(N\) such that \(q\not\equiv0\pmod{2^N}\).
The rational pairing shows that \(r=\rank L\) is positive.

Let \(\mathcal A=\phi_*B\) be a negative step returned to the \(Y_2\)
convention, and let \(\mathcal D=HB\) be a positive-return step.
By Theorems~\ref{thm:ordinary:units} and
\ref{thm:negative:unit}, and Lemma~\ref{lem:stack:finite-lattice},
both operators are automorphisms of \(L/2^NL\).  The chronological
four-interval pattern is one positive-return step followed by three
negative steps, so its operator is \(\mathcal A^3\mathcal D\).  Geometrically its
first two intervals and intervening positive bridge form the macro;
its last two intervals are single factors.  Proposition~
\ref{prop:stack:macro} proves that its directly counted clamped
macro has the same integral homology operation, with the already
fixed common sign convention.  For the finite-order argument it
would suffice that this direct operation is itself a mod-two unit.

Choose a common exponent \(E\) of
\(\GL_r(\Z/2^N\Z)\), enlarged to be divisible by eight.
Choose initial and terminal pad lengths \(p_-,p_+\) that are
multiples of \(E\) and meet the prescribed lower bounds.  Let
\(q_0\) be any sufficiently large positive multiple of \(E\)
and repeat the four-interval pattern \(q_0\) times.  Then
\[
 m=q_0,\qquad n=p_-+4q_0+p_+,
\]
and each pad and the whole repeated middle operation are the identity
on \(L/2^NL\).  Here \(q_0\), the repetition number, is distinct
from the fixed nonzero pairing \(q\).
Use literally periodic local block, orientation, and determinant
transition data within each set of repetitions.  Thus each operation
really is a power of one automorphism.  The data may depend on this
finite stack: the common exponent depends only on \(r,N\), so the
argument is unchanged if the automorphism changes when smaller
perturbations are required.  Continuation identifies the free
lattices, and the integral cap pairing remains \(q\), up to a
fixed sign, by the closed comparison in the proof of
Proposition~\ref{prop:stack:composition}.  A common sign of a map
causes no difficulty, since it too is an element of the same finite
group.

We check the degree congruence before making any metric choices.
The old nonzero cap pairing has an allowed bundle degree.  Rational
Mayer--Vietoris at its seam identifies its positive rank with
\(b^+(C_-)+b^+(C_+)=b_0\) and its determinant square with the sum
of the two cap squares.  Adjoining the outside trace halves adds only
negative directions with zero determinant evaluation.  Internal
negative cells also have \(c_0=0\), and a positive cell has
\(c_0^2=0\).  Thus the base determinant square of the closed stack
equals that for the old pairing, while its positive rank increases
by \(m\).  Equation~
\eqref{eq:stack:ordinary-dimension} therefore changes the required
quantity \(8\kappa\) by \(n+3m\).  Our choices make both
\(n\) and \(m\) multiples of eight, so this change is divisible
by eight.  The required \(c_2=\kappa+c_0^2/4\) is integral.
Every bit sector is also integral, by
\eqref{eq:stack:global-bit-squares}.  This is the residue restriction
on the allowed degrees; it has not been inferred from a formal
nonzero count.

The pad lengths and exceptional evaluations are now fixed.
Equation~\eqref{eq:stack:dirac-growth} becomes
\[
 n_D=\frac{q_0}{8}
 +\Theta_0-\frac{p_-+p_++z+3+3b_0}{8}.
\]
It is positive for all sufficiently large allowed \(q_0\).
Increasing the repetition number by \(E\) preserves every
congruence and increases \(n_D\) by \(E/8\).

For each such fixed finite problem choose the metric isolations and
small data as in Proposition~\ref{prop:stack:composition}, preserving
periodicity during the factor computation.  That proposition
identifies the actual count with the product.  Modulo \(2^N\)
this product pairs the final cap states exactly as the identity
operation, and hence
\[
 \Omega\equiv\pm q\not\equiv0\pmod{2^N}.
\]
It follows that \(\Omega\ne0\).  Subsequent sufficiently small
generic data for the one-component closed problem need no
periodicity and preserve this integer.  The same proposition gives
the compatibility with the metric tests and completes all assertions.
\end{proof}

With the stack integers and lifts fixed, the displayed indices bound
all charge ranges that will occur, including the auxiliary relative
cap calculations.  Table~\ref{tab:stack:choices} records the order of
the finite choices; its length and perturbation quantifiers are those
of Corollary~\ref{cor:estimates:choices}.

One deformation used in comparing the old cap closure with unperturbed
closed data deserves separate mention.  Turn off a gradient perturbation
while its seam remains long, preserving its gradient form throughout.
This retains the uniform charge bound used in the period exclusion.
\begin{table}[htbp]
\centering
\small
\begin{tabular}{@{}p{.04\textwidth}p{.32\textwidth}p{.56\textwidth}@{}}
\toprule
 & Choice & Dependence and use \\
\midrule
1 & Old caps, probes, and periods & The constants $C_W$ depend only on
the old cap data (Proposition~\ref{prop:estimates:cap}). \\
2 & Outside delays and odd exceptional evaluations & Chosen using the
old-cap bounds, before middle repetitions
(Lemma~\ref{lem:indices:end-exclusion}). \\
3 & Pads and repetition integers & Congruences, nonvanishing, and
$n_D>0$ determine a fixed finite problem
(Theorem~\ref{thm:stack:nonzero}). \\
4 & Tube scale $\epsilon$ & Chosen for this fixed finite problem. \\
5 & Regularizer $d$, finite rational isolations, and a positive gradient
tolerance & Isolation thresholds may depend on the whole finite
problem; the gradient tolerance precedes the external seam lengths. \\
6 & Periodic ordinary factor data, then finite $Y$ lengths & Data lie
below the chosen tolerances.  The $Y$ lengths may then be any sufficiently
large finite gluing lengths (Proposition~\ref{prop:stack:composition}). \\
7 & Further nongradient errors & Their allowed total squared cost may
depend on the selected finite $Y$ lengths. \\
8 & Small closed generic changes & These preserve the computed integer
$\Omega$ and need not preserve periodicity. \\
\bottomrule
\end{tabular}
\caption{Order of the finite choices.  Only the old-cap constants are
independent of later exterior choices.  The genuine $J$ and lens lengths,
distinct from the finite lengths listed here, may subsequently diverge
at parameter faces.}
\label{tab:stack:choices}
\end{table}

\subsection{The coupled spinor data}

The next sections couple the fixed spin-\(c\) determinant
\(l=\Lambda-c\) to the rank-two bundle of determinant \(c\).
Let $W^\pm$ be the rank-two spinor bundles of this spin-$c$ structure,
and let $E$ be the rank-two Hermitian bundle with $c_1(E)=c$.
A configuration consists of a unitary connection $a$ on $E$, with
its determinant connection fixed, and a coupled spinor
$\Phi\in\Gamma(W^+\otimes E)$.  The coupled Dirac operator is
$D_a:\Gamma(W^+\otimes E)\longrightarrow\Gamma(W^-\otimes E)$.
In the equations of Section~\ref{sec:indices}, $P_D$ is a section of
$W^-\otimes E$, while $P_+$ has the same target as
$\rho(F_a^{0,+})$.  These denote the Dirac and curvature perturbations.
The phase circle multiplies $\Phi$; on the determinant-one gauge
quotient its ineffective subgroup at a free orbit is $\{1,-1\}$.
Definition~\ref{def:indices:types} gives the full field-type
classification.  Table~\ref{tab:gluing:orientations} records how the
ordinary determinant, framing, and sign conventions continue into
the final spinor count.

At a phase-fixed configuration the spinor is either zero, or lies
in one line of a parallel splitting.  If \(v\) is the ordered
summand difference, with the active line first in the latter case,
its spin-\(c\) line class is \(K=\Lambda+v\).  For virtual
\(-4\)-cap fillings use the split references above, choosing the
order at the nonextendible center so that \(|K S_i|\le4\);
this is possible because \(|\Lambda S_i|=2\), and is automatic
at the trace-zero state.  These conventions fix the interface with
Section~\ref{sec:indices}.  The symbol \(P\) for a positive bridge
cell in this section is distinct from the later component type
\(P\), which denotes a field not fixed by phase.
